% archon:covers Smale5/UpstreamBridge.lean Smale5/Orders.lean Smale5/Global.lean Smale5/PrincipalDivisor.lean
\chapter{Checked implications with explicit premises}

\begin{lemma}[Order arithmetic]
\label{bp:order-arithmetic}
\uses{bp:products}
\lean{Smale5.sum_fourOrders, Smale5.exists_negative_fourOrder, Smale5.diagram9_totalOrder_pos, Smale5.diagram9_unbounded_fourProducts}
\leanok
The sum of the ten four-product orders is four times the total distance
order. A negative total order forces a negative four-product order.
The supplied diagram-9 order pattern gives positive total order and negative
orders for the specified four-products. This verifies consequences of the
pattern, not the analytic classification that supplies the pattern.
\end{lemma}

\begin{lemma}[Degree-zero arithmetic]
\label{bp:degree-zero-arithmetic}
\lean{Smale5.exists_negative_of_sum_zero_of_positive, Smale5.zero_has_pole_of_degree_zero}
\leanok
For an order function supported on a finite set with total sum zero, a
positive value forces a negative value. Finite support and degree zero
are explicit hypotheses; no geometric principal-divisor theorem is assumed
as a new Lean axiom.
\end{lemma}

\begin{theorem}[Conditional global contradiction]
\label{bp:conditional-global}
\uses{bp:order-arithmetic, bp:degree-zero-arithmetic}
\lean{Smale5.globalOrderContradiction, Smale5.globalOrderContradiction_from_degreeZero, Smale5.finite_of_order_reduction}
\leanok
Suppose an initial four-product pole exists, every four-product pole forces
a positive total order, and a positive total order forces a negative one.
The order arithmetic gives a contradiction. The version using finite support
and degree zero proves the second implication from those premises.
The general finiteness reduction additionally assumes that an infinite set
produces such order data. The required data have not yet been constructed
for the five-body configuration space.
\end{theorem}

\begin{theorem}[Adapter to the published upstream main theorem]
\label{bp:upstream-adapter}
\uses{bp:configuration, bp:distance-fibre}
\lean{Smale5.positive_normalized_isCC, Smale5.finite_positive_normalized_similar, Smale5.finite_positive_normalized_of_upstream, Smale5.fiveBodyFiniteness_of_upstream}
\leanok
Positive normalized configurations satisfy the genuine Newtonian central
configuration equations. Similar normalized configurations have identical
mutual distances; each similarity class therefore has a finite normalized
fibre. The literal upstream all-positive-mass finiteness theorem implies the
project's exact FiveBodyFiniteness target. That upstream theorem is an explicit
argument here, not an axiom or a claim that the expensive external project
has already been rebuilt locally. The separate integration source applies
the published CC5Full main theorem.
\end{theorem}
